\begin{tcolorbox}[colback=red!3!white,colframe=red!50!black,boxrule=0.8pt,title={Correction notes},fonttitle=\bfseries]
\textbf{Corrections from the calculation check (30~Sept.~2026).} The following places are corrected or annotated in the text; each carries a dated note with the previous value:
\begin{itemize}
\item Doc.~154: Gaussian curvature $\theta \approx \pi$ ``outside'' $\to$ ``inside''; $r_{\min} \approx 21\,\ell_P \approx 3.4\times10^{-34}$\,m $\to$ $19.6\,\ell_P \approx 3.2\times10^{-34}$\,m; ``variational calculus shows'' (torus most stable form) $\to$ conjecture [S].
\end{itemize}
\end{tcolorbox}

\section*{Abstract}
		This paper examines the parallel between brain folding (cortical gyri) and the 4D torsional structure of T0 theory. The metaphor can be stated mathematically at several points; where it goes beyond a structural analogy, the text says so. Both systems solve the same fundamental problem: \textbf{How does one pack maximum surface area/information into minimum volume without singularities?} The analysis describes nine parallels: (1) \textbf{Fractal self-similarity} across many scales. (2) \textbf{Surface maximization} with volume minimization. (3) \textbf{Deep furrows = high density}: Sulci $\leftrightarrow$ mass concentrations. (4) \textbf{Singularity avoidance} through minimum curvature radius. (5) \textbf{Static structure, dynamic flows}: Material static, information dynamic. (6) \textbf{Hierarchical information processing} across levels. (7) \textbf{Topological invariants}: Genus = 1 for both. (8) \textbf{Energy efficiency} through geometric optimization. (9) \textbf{Asymmetry as function}: Left vs. right hemisphere $\leftrightarrow$ cosmic dipoles. The parallels suggest that both systems use a similar geometric solution for information storage and processing; they are to be read as a structural analogy, not as proof of a common mechanism.

	\medskip\noindent\textbf{Note on versions.} This English version of the document is more extensive than the German one. Earlier revisions shortened or changed parts of the German version. Where the two differ, this more extensive English version applies.


	\medskip\noindent\textbf{Status markers.} Statements marked \textbf{[S]} are \emph{speculative}: a hint or conjecture that has not yet been derived or numerically checked in the corpus. Here this mainly concerns the cosmological statements; the current corpus deliberately sets the cosmic sector aside, since the Standard Model does not derive it either (P39).

	
	
	
	\section{Introduction: The Image}
	
	\subsection{The Metaphor}
	
	In FFGF/T0 theory, the universe is described as:
	
	\begin{center}
		\textbf{"A huge, fractally folded brain"}
	\end{center}
	
	where the **deep folds** (sulci) correspond to regions of highest mass and energy density.
	
	\subsection{Why is this metaphor fitting?}
	
	\begin{important}[Central Observation]
		The human brain and the universe in T0 theory solve **the same fundamental optimization problem**:
		
		\vspace{0.3cm}
		
		\textbf{How does one maximize surface area (information, density) in minimum volume without creating singularities (collapse)?}
		
		\vspace{0.3cm}
		
		The answer in both cases: \textbf{fractal folding}.
	\end{important}
	
	\section{Nine Parallels}
	
	\subsection{Parallel 1: Fractal Self-Similarity}
	
	\subsubsection{Brain}
	
	The human cortex shows fractal structure:
	\begin{itemize}
		\item \textbf{Large furrows} (primary sulci): 1--2 cm deep
		\item \textbf{Medium convolutions} (secondary sulci): 0.5--1 cm
		\item \textbf{Small folds} (tertiary sulci): 0.1--0.5 cm
		\item \textbf{Microcolumns}: 30--50 $\mu$m
	\end{itemize}
	
	Each large fold contains smaller folds following the same principle.
	
	\textbf{Fractal dimension of cortex}: $D_{\text{cortex}} \approx 2.7 - 2.8$
	
	\subsubsection{T0 Universe}
	
	The torus structure scales self-similarly over **60+ orders of magnitude**:
	
	\begin{table}[H]
		\centering
		\small
		\adjustbox{max width=\textwidth}{%
\begin{tabular}{lll}
			\toprule
			\textbf{Scale} & \textbf{R (major radius)} & \textbf{System} \\
			\midrule
			Sub-Planck & $\sim 10^{-39}$ m & Fundamental granulation \\
			Particles & $\sim 10^{-15}$ m & Protons, leptons \\
			Atoms & $\sim 10^{-10}$ m & Electron shells \\
			Planets & $\sim 10^{6}$ m & Magnetic field torus \\
			Stars & $\sim 10^{9}$ m & Convection currents \\
			Galaxies & $\sim 10^{20}$ m & Spiral arms \\
			Cosmic web & $\sim 10^{24}$ m & Filaments \\
			\bottomrule
		\end{tabular}%
}
		\caption{Self-similar torus structures across scales}
	\end{table}
	
	\textbf{Fractal dimension}: $D_f = 3 - \xi \approx 2.9998666$
	
	\begin{keyresult}[First Parallel]
		Both systems show **fractal self-similarity**: Each large structure contains smaller versions following the same geometric principle.
		
		\textbf{Mathematically}: Both fractal dimensions lie between 2 and 3, clearly below 3 for the cortex and very close to 3 for the T0 universe.
		\begin{itemize}
			\item Cortex: $D \approx 2.75$
			\item Universe: $D \approx 2.9998666$
		\end{itemize}
	\end{keyresult}
	
	\subsection{Parallel 2: Surface Maximization}
	
	\subsubsection{Brain}
	
	\textbf{Problem}: How to pack $\sim 16$ billion neurons into a skull of $\sim 1.3$ liters?
	
	\textbf{Solution}: Folding increases surface area.
	
	\begin{align}
		\text{Smooth sphere} &\to A = 4\pi r^2 \approx 600\,\text{cm}^2 \\
		\text{Folded cortex} &\to A \approx 2400\,\text{cm}^2
	\end{align}
	
	\textbf{Factor 4 more surface area} through folding at same volume.
	
	\subsubsection{T0 Universe}
	
	\textbf{Problem}: How to pack maximum energy density into minimum space without singularities?
	
	\textbf{Solution}: Torus folding.
	
	For a torus:
	\begin{align}
		\text{Surface area} &: A = 4\pi^2 R r \\
		\text{Volume} &: V = 2\pi^2 R r^2 \\
		\text{Ratio} &: \frac{A}{V} = \frac{2}{r}
	\end{align}
	
	The smaller $r$ (tube radius), the **greater the surface area per volume**.
	
	\textbf{Limit}: $r_{\min} \approx 21 \ell_P$ prevents singularity.
	
	\begin{keyresult}[Second Parallel]
		Both systems maximize surface area at minimum volume:
		
		\begin{itemize}
			\item \textbf{Brain}: Maximum neuronal surface area
			\item \textbf{Universe}: Maximum energy density surface area
		\end{itemize}
		
		\textbf{Both avoid singularities}:
		\begin{itemize}
			\item Cortex: Minimum sulcus depth $\sim 1$ mm (blood supply)
			\item Universe: Minimum radius $r_{\min} = 21 \ell_P$
		\end{itemize}
	\end{keyresult}
	
	\subsection{Parallel 3: Deep Furrows = High Density}
	
	\subsubsection{Brain}
	
	The **deepest sulci** (furrows) of the brain contain the **densest neuronal connections**:
	
	\begin{itemize}
		\item \textbf{Lateral fissure} (Sylvian fissure): Separation frontal/temporal lobes
		\begin{itemize}
			\item $\to$ Language centers (Broca, Wernicke)
			\item $\to$ high cognitive density
		\end{itemize}
		\item \textbf{Central sulcus}: Motor/sensory cortex
		\begin{itemize}
			\item $\to$ Direct body control
			\item $\to$ Maximum information density
		\end{itemize}
	\end{itemize}
	
	\textbf{Principle}: Deep folds $\leftrightarrow$ high functional importance
	
	\subsubsection{T0 Universe}
	
	The **deepest folds** of the torus geometry (regions with negative Gaussian curvature) correspond to **highest mass densities**:
	
	Gaussian curvature of torus:
	\begin{equation}
		K(\theta) = \frac{\cos\theta}{r(R + r\cos\theta)}
	\end{equation}
	
	\textbf{Inside} ($\theta \approx \pi$): $K < 0$ \textit{(Corrected on 30~Sept.~2026: previously ``Outside'' -- with $K(\theta) = \cos\theta/(r(R + r\cos\theta))$, $\theta = \pi$ lies at distance $R - r$ from the axis, i.e.\ inside; outside ($\theta = 0$), $K = 1/(r(R+r)) > 0$.)} $\to$ \textbf{Negative curvature}
	
	Here T0 theory places \textbf{[S]}:
	\begin{itemize}
		\item Galaxy cores
		\item Supermassive black holes
		\item Supercluster nodes
		\item Filament intersection points
	\end{itemize}
	
	\begin{keyresult}[Third Parallel]
		\textbf{Deep furrows = High density}
		
		\vspace{0.3cm}
		
		\begin{center}
			\adjustbox{max width=\textwidth}{%
\begin{tabular}{p{4cm}|p{4cm}}
				\toprule
				\textbf{Brain} & \textbf{Universe (T0)} \\
				\midrule
				Deepest sulci & Negative curvature ($K<0$) \\
				$\downarrow$ & $\downarrow$ \\
				Densest neuronal connections & Highest mass density \\
				$\downarrow$ & $\downarrow$ \\
				Maximum information & Maximum energy \\
				\bottomrule
			\end{tabular}%
}
		\end{center}
	\end{keyresult}
	
	\subsection{Parallel 4: Singularity Avoidance}
	
	\subsubsection{Brain}
	
	The cortex cannot fold **arbitrarily deep**:
	
	\textbf{Limitations}:
	\begin{enumerate}
		\item \textbf{Blood supply}: Deep sulci need capillaries
		\item \textbf{Mechanical stability}: Too thin walls collapse
		\item \textbf{Minimum thickness}: $\sim 1.5 - 4$ mm (gray/white matter)
	\end{enumerate}
	
	$\Rightarrow$ Minimum curvature radii prevent "singularities"
	
	\subsubsection{T0 Universe}
	
	The fractal dimension $D_f = 3 - \xi$ prevents collapse:
	
	In perfect 3D space ($D=3$): Torus could shrink to $r \to 0$ (singularity).
	
	With $D_f = 3-\xi$: Minimum tube radius
	\begin{equation}
		r_{\min} \propto \frac{\ell_P}{\xi^{1/3}} \approx 19.6 \times \ell_P \approx 3.2 \times 10^{-34}\,\text{m}
	\end{equation}
	\textit{(Corrected on 30~Sept.~2026: previously $21 \times \ell_P \approx 3.4 \times 10^{-34}$\,m -- $\xi^{-1/3} = (7500)^{1/3} = 19.57$, $19.57 \cdot 1.616 \times 10^{-35}\,\text{m} = 3.16 \times 10^{-34}$\,m.)}
	
	\textbf{Meaning}: In the T0 description, the fractal structure of space itself prevents singularities.
	
	\begin{keyresult}[Fourth Parallel]
		\textbf{Both systems avoid singularities} through natural minimum curvature radii:
		
		\begin{itemize}
			\item \textbf{Brain}: $r_{\min} \sim 1$ mm (biological)
			\item \textbf{Universe}: $r_{\min} \sim 21 \ell_P$ (geometrical)
		\end{itemize}
		
		The folding increases surface area **without collapsing into singularities**.
	\end{keyresult}
	
	\subsection{Parallel 5: Static + Dynamic}
	
	\subsubsection{Brain}
	
	\textbf{Structure}: Materially **static**
	\begin{itemize}
		\item Neurons don't move
		\item Cortex architecture is fixed
		\item Anatomy remains constant
	\end{itemize}
	
	\textbf{Function}: Electrically **dynamic**
	\begin{itemize}
		\item Action potentials propagate
		\item Synapses fire
		\item Information flows
	\end{itemize}
	
	\subsubsection{T0 Universe}
	
	\textbf{Structure} \textbf{[S]}: In this reading the universe is **static**
	\begin{itemize}
		\item No Big Bang
		\item No cosmic expansion
		\item 4D torsion crystal is timeless
	\end{itemize}
	
	\textbf{Dynamics}: Energy flows are **dynamic**
	\begin{itemize}
		\item Photons propagate
		\item Torsion waves travel
		\item Energy circulates in torus
	\end{itemize}
	
	\textbf{Redshift} \textbf{[S]}: In this reading it arises not from expansion, but from:
	\begin{equation}
		z \approx \xi \cdot \ln\left(\frac{d}{\ell_P}\right)
	\end{equation}
	Fractal energy loss along the folds. These cosmological statements are speculative; the cosmic sector is deliberately set aside in the current corpus (P39).
	
	\begin{keyresult}[Fifth Parallel]
		\textbf{Static base structure, dynamic flows}:
		
		\vspace{0.3cm}
		
		\begin{center}
			\adjustbox{max width=\textwidth}{%
\begin{tabular}{lcc}
				\toprule
				& \textbf{Brain} & \textbf{Universe (T0)} \\
				\midrule
				Material/Structure & Static & Static \\
				Information/Energy & Dynamic & Dynamic \\
				Surface/Space & Folded & Folded (torus) \\
				\bottomrule
			\end{tabular}%
}
		\end{center}
	\end{keyresult}
	
	\subsection{Parallel 6: Hierarchical Processing}
	
	\subsubsection{Brain}
	
	Neuronal information processing is **hierarchical**:
	
	\begin{enumerate}
		\item \textbf{Level 1}: Receptors (retina, cochlea)
		\item \textbf{Level 2}: Primary sensory areas (V1, A1)
		\item \textbf{Level 3}: Secondary areas (V2, V4)
		\item \textbf{Level 4}: Association cortex
		\item \textbf{Level 5}: Prefrontal cortex (executive function)
	\end{enumerate}
	
	Each level extracts more abstract features.
	
	\subsubsection{T0 Universe}
	
	Torsion structures are nested across scales:
	
	\begin{enumerate}
		\item \textbf{Sub-Planck}: $\Lambda_0 \sim 10^{-39}$ m -- Fundamental granulation
		\item \textbf{Planck}: $\ell_P \sim 10^{-35}$ m -- Quantum gravity
		\item \textbf{Particles}: $\sim 10^{-15}$ m -- Protons, leptons
		\item \textbf{Atoms}: $\sim 10^{-10}$ m -- Electron shells
		\item \textbf{Stars}: $\sim 10^9$ m -- Convection torus
		\item \textbf{Galaxies}: $\sim 10^{20}$ m -- Spiral arms
		\item \textbf{Cosmic}: $\sim 10^{24}$ m -- Filament network
	\end{enumerate}
	
	Each scale is described as a torus, **embedded in larger tori**.
	
	\begin{keyresult}[Sixth Parallel]
		\textbf{Hierarchical information processing}:
		
		\begin{itemize}
			\item \textbf{Brain}: Neural networks on different abstraction levels
			\item \textbf{Universe}: Nested torus vortices from Planck to Hubble
		\end{itemize}
		
		Both are **fractally layered**.
	\end{keyresult}
	
	\subsection{Parallel 7: Topological Invariance}
	
	\subsubsection{Brain}
	
	Topologically, the cortex can be regarded as a **torus**.
	
	\textbf{Why?}
	\begin{itemize}
		\item Cerebral hemispheres are connected by the **corpus callosum**
		\item The ventricular system forms a **central hole**
		\item Genus = 1 (one hole)
	\end{itemize}
	
	Mathematically: In this simplification, the folded cortex can be continuously deformed into a torus.
	
	\subsubsection{T0 Universe}
	
	The fundamental structure is a **4D torus**:
	
	\begin{equation}
		\mathcal{M} = \mathbb{R}^3 \times S^1_{\text{comp}}
	\end{equation}
	
	\textbf{Properties}:
	\begin{itemize}
		\item 3 spatial + 1 compact dimension
		\item Genus = 1 (one hole)
		\item Poloidal + toroidal circulation
	\end{itemize}
	
	\begin{keyresult}[Seventh Parallel]
		\textbf{Both have the same topology: Torus (Genus = 1)}
		
		At the level of topology this is more than a metaphor: both have the same genus.
		
		\begin{itemize}
			\item Cortex: Topologically equivalent to torus
			\item Universe: Fundamental 4D torus
		\end{itemize}
		
		The topology is **invariant** under folding.
	\end{keyresult}
	
	\subsection{Parallel 8: Energy Efficiency}
	
	\subsubsection{Brain}
	
	The brain is **very energy efficient**:
	
	\begin{itemize}
		\item Power: $\sim 20$ Watts
		\item Operations: $\sim 10^{16}$ synapses/second
		\item Efficiency: $\sim 10^{-15}$ Joules per operation
	\end{itemize}
	
	\textbf{Reason}: Folding shortens wiring (axons) while keeping connectivity high.
	
	\textbf{Principle}: Minimize
	\begin{equation}
		E_{\text{total}} = E_{\text{wiring}} + E_{\text{volume}}
	\end{equation}
	
	$\Rightarrow$ Solution: folded surface.
	
	\subsubsection{T0 Universe}
	
	The torus minimizes energy for given topology:
	
	\begin{equation}
		E_{\text{total}} = E_{\text{surface}} + E_{\text{curvature}} + E_{\text{rotation}}
	\end{equation}
	
	\textbf{Conjecture \textbf{[S]}}: For constant flux and angular momentum, the torus is the **most stable form**. \textit{(Corrected on 30~Sept.~2026: previously ``Variational calculus shows'' -- the variational calculation has not been carried out, as Doc.~145 explicitly states; the statement is speculative [S].)}
	
	The fractal dimension $D_f = 3-\xi$ means:
	\begin{itemize}
		\item Energy experiences "resistance" when flowing
		\item Torus is the path of **least resistance**
	\end{itemize}
	
	\begin{keyresult}[Eighth Parallel]
		\textbf{Both systems optimize energy}:
		
		\begin{itemize}
			\item \textbf{Brain}: Minimum wiring, maximum function
			\item \textbf{Universe}: Minimum energy, maximum stability
		\end{itemize}
		
		The folding is the **solution to a variational problem**.
	\end{keyresult}
	
	\subsection{Parallel 9: Asymmetry as Function}
	
	\subsubsection{Brain}
	
	The brain is **asymmetric**:
	
	\begin{itemize}
		\item \textbf{Left hemisphere}: Language, logic, sequential
		\item \textbf{Right hemisphere}: Spatial, holistic, parallel
	\end{itemize}
	
	This asymmetry is **functional**, not a defect.
	
	\textbf{Folding pattern}: Left and right different
	\begin{itemize}
		\item Left Sylvian fissure: Deeper (language center)
		\item Right parietal lobe: Larger (spatiality)
	\end{itemize}
	
	\subsubsection{T0 Universe}
	
	The universe shows **asymmetries** that T0 theory interprets as intrinsic \textbf{[S]}:
	
	\begin{itemize}
		\item \textbf{CMB dipole}: Preference direction in cosmic microwave background
		\item \textbf{Cosmic flows}: Large-scale movements
		\item \textbf{Two-dipole model}: Fundamental asymmetry of the "global fold"
	\end{itemize}
	
	In T0 theory this asymmetry is interpreted as structural rather than as a disturbance (speculative, cosmic sector, P39).
	
	It arises from **pentagonal symmetry breaking** by the golden ratio $\varphi$:
	\begin{equation}
		\xi = \frac{4}{30000} \quad \text{with factor } 5\varphi \text{ in the structure}
	\end{equation}
	
	\begin{keyresult}[Ninth Parallel]
		\textbf{Asymmetry is functional}:
		
		\vspace{0.3cm}
		
		\begin{center}
			\adjustbox{max width=\textwidth}{%
\begin{tabular}{p{4cm}|p{4cm}}
				\toprule
				\textbf{Brain} & \textbf{Universe (T0)} \\
				\midrule
				Left vs. right hemisphere & CMB dipole, cosmic flows \\
				Functional specialization & Global asymmetry of fold \\
				Emerges from development & Emerges from $\varphi$-breaking [S] \\
				\bottomrule
			\end{tabular}%
}
		\end{center}
	\end{keyresult}
	
	\section{How far does the metaphor carry?}
	
	\subsection{Mathematical Precision}
	
	Some of the parallels can be stated **quantitatively**:
	
	\begin{table}[H]
		\centering
		\small
		\adjustbox{max width=\textwidth}{%
\begin{tabular}{lcc}
			\toprule
			\textbf{Property} & \textbf{Brain} & \textbf{Universe (T0)} \\
			\midrule
			Fractal dimension & $D \approx 2.75$ & $D_f = 3-\xi \approx 2.9998666$ \\
			Topological genus & 1 (torus) & 1 (4D torus) \\
			Surface gain & $\times 4$ & $\propto 1/r_{\min}$ \\
			Minimum radius & $\sim 1$ mm & $21 \ell_P$ \\
			Hierarchy levels & $\sim 5-6$ & $>60$ \\
			\bottomrule
		\end{tabular}%
}
		\caption{Quantitative parallels}
	\end{table}
	
	\subsection{Universal Optimization Principle}
	
	Both solve a comparable problem through **a similar geometric strategy**:
	
	\begin{center}
		
		\textbf{Maximize $\frac{\text{Surface (Information)}}{\text{Volume (Space)}}$}
		
		\normalsize
		under the constraint:
		
		
		\textbf{No singularities.}
	\end{center}
	
	\subsection{Information is Geometry}
	
	The interpretation that follows:
	
	\begin{revolutionary}[Information = Geometry]
		\textbf{In both systems, information is geometrically encoded.}
		
		\vspace{0.3cm}
		
		\textbf{Brain}:
		\begin{itemize}
			\item Neuronal information $\leftrightarrow$ folding structure
			\item More surface = more synapses = more information
		\end{itemize}
		
		\textbf{Universe}:
		\begin{itemize}
			\item Physical information $\leftrightarrow$ torsion structure
			\item More windings = more energy = more information
		\end{itemize}
		
		\vspace{0.3cm}
		
		The metaphor suggests: **geometry carries information**.
	\end{revolutionary}
	
	\section{The Narrative Power}
	
	\subsection{Why brain instead of other metaphors?}
	
	There are many folded systems (paper, fabric, intestine, ...). Why is the **brain** so fitting?
	
	\begin{philosophical}[Why brain?]
		\textbf{1. Consciousness and cosmos}:
		
		The brain is the most complex known object in the universe. The metaphor suggests: The universe itself might have a form of "consciousness" -- not in an anthropomorphic sense, but as a **self-organizing information system**.
		
		\vspace{0.3cm}
		
		\textbf{2. Micro-macro unity}:
		
		The smallest conscious system (brain, $\sim$1 kg) and the largest system (universe, $\sim 10^{53}$ kg) follow, in this reading, **similar geometric principles**.
		
		This corresponds to a core statement of T0 theory: **self-similarity over 60 orders of magnitude**.
		
		\vspace{0.3cm}
		
		\textbf{3. Emergence and complexity}:
		
		From simple folding rules (torus geometry) emerges great complexity:
		\begin{itemize}
			\item Brain: $\sim 86$ billion neurons, $\sim 10^{14}$ synapses
			\item Universe: $\sim 10^{80}$ particles, cosmic web
		\end{itemize}
		
		Both are **more than the sum of their parts**.
	\end{philosophical}
	
	\subsection{The holographic principle}
	
	The brain folding metaphor connects with the **holographic principle**:
	
	\begin{keyresult}[Holography]
		\textbf{Holographic principle}: The information of a volume is encoded on its surface.
		
		\vspace{0.3cm}
		
		\textbf{Brain}: The $\sim 2$ mm thin cortex **surface** contains all cognitive information -- the underlying volume (white matter) serves mainly as wiring.
		
		\vspace{0.3cm}
		
		\textbf{Universe (T0)}: The torsion **surface** (4D hypersurface) encodes all physical information -- the "volume" is emergent \textbf{[S]}.
		
		\vspace{0.3cm}
		
		Folding increases surface $\Rightarrow$ increases information capacity.
	\end{keyresult}
	
	\section{Summary: Nine Parallels}
	
	\begin{table}[H]
		\centering
		\footnotesize
		\resizebox{\textwidth}{!}{%
		\begin{tabular}{clp{4cm}p{4cm}}
			\toprule
			\textbf{No.} & \textbf{Parallel} & \textbf{Brain} & \textbf{Universe (T0)} \\
			\midrule
			1 & Fractal self-similarity & Sulci at all scales & Torus structures 60+ orders of magnitude \\
			2 & Surface maximization & $\times 4$ through folding & $\propto 1/r_{\min}$ \\
			3 & Deep furrows = density & Neuronal density in sulci & Mass density at $K<0$ \\
			4 & Singularity avoidance & $r_{\min} \sim 1$ mm & $r_{\min} = 21 \ell_P$ \\
			5 & Static + dynamic & Material static, electrical dynamic & Structure static, energy dynamic \\
			6 & Hierarchical processing & 5-6 cortical levels & 7+ scale levels \\
			7 & Topology: torus & Genus = 1 & 4D torus \\
			8 & Energy efficiency & Minimum wiring & Minimum energy \\
			9 & Asymmetry as function & Left vs. right & CMB dipole [S] \\
			\bottomrule
		\end{tabular}}
		\caption{The nine parallels}
	\end{table}
	
	\section{Conclusion}
	
	\begin{keyresult}[Why does the metaphor fit?]
		The brain folding metaphor fits for the following reasons:
		
		\vspace{0.3cm}
		
		\textbf{1. Mathematical common ground}:
		Both have fractal dimensions in the range $D \approx 2.7 - 3.0$ and can be described by torus topology (genus = 1).
		
		\vspace{0.3cm}
		
		\textbf{2. Same optimization problem}:
		Both maximize surface area/information at minimum volume without singularities.
		
		\vspace{0.3cm}
		
		\textbf{3. Self-similarity}:
		Both show fractal hierarchy across many scales.
		
		\vspace{0.3cm}
		
		\textbf{4. Information = geometry}:
		Both encode information in folded surface.
		
		\vspace{0.3cm}
		
		\textbf{5. Narrative depth}:
		The metaphor connects the smallest conscious system (brain) with the largest system (universe) and raises the question whether **consciousness and cosmos are geometrically related** \textbf{[S]}.
		
		\vspace{0.3cm}
		
		The parallels point to a **common geometric solution** for information storage and processing; whether a common principle lies behind it remains an open question.
	\end{keyresult}
	
	\subsection{Summary interpretation}
	
	\begin{revolutionary}[Interpretation]
		
		\begin{center}
			\textbf{The universe doesn't think like a brain --}
			
			\vspace{0.3cm}
			
			\textbf{The brain is folded like the universe.}
		\end{center}
		
		\normalsize
		
		\vspace{0.3cm}
		
		Both can be described with the same geometric logic:
		
		\begin{equation}
			\boxed{\max\left(\frac{\text{Surface}}{\text{Volume}}\right) \text{ with } r \geq r_{\min}}
		\end{equation}
		
		The solution in both cases: \textbf{fractal folding in torus topology}.
	\end{revolutionary}
	